% chapter-1.tex
% Chapter 1 — Introduction
%
% Citations use the citation keys of the project Zotero library; the
% corresponding entries were exported to references.bib. Claims that would
% require a source not currently available in the library are marked with
% % TODO comments and should be resolved before submission.

\chapter{Introduction}
\label{chap:introduction}

\section{Background}
\label{sec:background}

Modern robotic and autonomous systems increasingly rely on computer vision
to perceive and interpret their surroundings. Tasks such as object
detection, classification, and tracking are dominated by deep neural
networks and are computationally demanding; in many applications they must
be executed continuously, in real time or near real time, on hardware
located on board the robot itself. The proliferation of mobile and
Internet-of-Things devices has moved the production of data towards the
network edge, creating pressure to run artificial-intelligence workloads
next to the sensors that generate the data rather than in remote data
centers---a paradigm consolidated under the term \emph{edge intelligence}
\cite{zhou2019, deng2020}. Edge deep learning integrates computational
resources and data sources closely, enabling real-time decision making
attuned to the environment \cite{xu2025}. Offloading perception to remote
servers is often impractical because of latency, bandwidth, and connectivity
constraints, and, in some applications, data-privacy considerations.
% TODO: the privacy motivation is not explicitly stated in the sources
% currently available in the library; find a dedicated source or remove
% the clause.
As a consequence, computer-vision workloads are commonly executed on
embedded computing platforms that must balance computational capability
against strict power, thermal, and cost budgets \cite{zhou2019, xu2025}.
The range of applications is broad and growing: autonomous drones for
wildlife monitoring \cite{lundquist2026}, edge-based defect detection in
agriculture and industry \cite{muchtar2025, machado2025}, smart-city
services \cite{woo2025}, and medical image analysis \cite{aalishah2025}.

Among embedded platforms, the NVIDIA Jetson family is representative of the
computing resources available to robots and edge artificial-intelligence
systems. These modules integrate Arm CPU cores with NVIDIA GPU accelerators
in a power-efficient form factor and are widely used for on-board computer
vision and deep-learning inference \cite{scalcon2024}. This thesis uses the
NVIDIA Jetson Orin Nano as its single experimental platform. Recent studies
indicate that the device is a credible host for exactly the class of
workload studied here: Scalcon et al.\ evaluated the Orin family for
AI-powered video monitoring and found the 15~W configuration of the Orin
Nano suitable for surveillance applications, sustaining on average 13~FPS
for YOLOv4 with DeepStream \cite{scalcon2024}, while Alqahtani et al.,
benchmarking object-detection models across heterogeneous edge devices,
concluded that the Orin Nano offers the most favorable overall balance
of accuracy, latency, and energy across most model families
\cite{alqahtani2026}.
Complementary work characterizes how neural-network inference on the Orin
Nano depends on the software stack---framework, numerical precision, and
language \cite{szegedi2026}---and on system-level energy and efficiency
trade-offs \cite{zequn2026}. Notably for this thesis, Szegedi et al.\
report that for object detection, preprocessing and postprocessing
represent a substantial portion of total runtime, especially at higher
input resolutions \cite{szegedi2026}: a computer-vision application is more
than its inference engine. It comprises a chain of stages---image
acquisition, preprocessing, inference, and postprocessing---and how those
stages are connected, within a single process or across communicating
processes, is an architectural decision with performance consequences of
its own. That decision is the subject of this thesis.

On the software side, robotic systems have long embraced modular
architectures. The Robot Operating System (ROS), introduced in 2009, is not
an operating system in the traditional sense of process management and
scheduling, but a structured communications layer above the host operating
system \cite{quigley2009}. Its publication-based model and reusable
components demonstrably expedited robotics research \cite{macenski2022}.
ROS~2, its second generation, was redesigned from the ground up to meet
requirements that ROS~1 did not anticipate, including support for real-time
systems and fully distributed, multi-robot deployments
\cite{macenski2022, portugal2025}. Architecturally, ROS~2 replaced the
central master of ROS~1 with a decentralized communication layer built on
the Data Distribution Service (DDS) standard, which contributes discovery,
configurable quality-of-service (QoS) policies, and transport primitives
suited to distributed real-time embedded systems
\cite{macenski2022, maruyama2016, gambo2025}. DDS-based middleware has
become widespread in robotics \cite{gambo2025}, and ROS~2 is currently used
across domains as diverse as autonomous vehicles, healthcare, agriculture,
aerospace, and logistics \cite{al-batati2026}, with its real-time
properties the subject of active research \cite{casini2026}. With
long-term support for ROS~1 ending in May 2025, the migration of the
community to ROS~2 is no longer optional \cite{portugal2025}: how to
structure applications in ROS~2 is a present and pressing design question
for essentially every team in this ecosystem.

In ROS~2, applications are decomposed into \emph{nodes}: independent
processes that exchange typed messages over logical channels called
\emph{topics}. This node-based model is the source of ROS~2's celebrated
modularity---separation of concerns, independent development and testing,
and reuse of components across systems \cite{macenski2022}---but it also
places inter-process communication on the critical path of the
application. Every message that crosses a node boundary may incur message
handling, serialization, memory copies, scheduler interaction, and
transport processing. This overhead is not hypothetical. Maruyama et al.\
observed, already in 2016, that ROS~2 must convert data for DDS and
abstract DDS from its users, and that this incurs additional overhead
\cite{maruyama2016}. Kronauer et al.\ profiled multi-node ROS~2 systems
and found that the ROS~2 layers add up to 50\% latency overhead relative
to raw DDS communication for small messages, with the magnitude depending
on the DDS implementation, payload size, and configuration
\cite{kronauer2021}. Kamiya et al.\ traced latency through the internal
layers of ROS~2 communication and attributed more than half of the total
latency in default socket-based communication to the DDS transport
segment; replacing socket transport with shared-memory communication
reduced overall latency to less than half \cite{kamiya2025}. Luo et al.\
showed that even intra-process communication, often assumed to be
negligible, exhibits costs that grow with message size and subscriber
count \cite{luo2025}. On embedded hardware these effects are amplified:
on a Raspberry Pi~4, Kronauer et al.\ measured the loss of roughly 30\%
of messages for 500~KB payloads and an extreme, non-linear increase in
latency once pipelines exceeded about 13--17 nodes at publishing
frequencies above 50~Hz \cite{kronauer2021}. The resource cost is
equally tangible: composing ROS~2 nodes into fewer processes saved 28\%
CPU and 33\% RAM for a complete navigation system executing on an ARM
processor, and Macenski et al.\ report that processing pipelines carrying
messages larger than 1~MB---with images and point clouds given as the
canonical examples---are virtually impossible without such composition
\cite{macenski2023}. Video streams are among the largest and
most frequent payloads in robotic systems, since autonomous platforms
process large volumes of sensor data at high rates \cite{kamiya2025}.

The practical significance of these overheads depends on the relation
between communication costs and the computing budget of the platform. On a
workstation-class machine, the cost of moving a frame between processes
may be small compared with the cost of computing on it. On a
resource-constrained embedded module operating under a power envelope of a
few watts, the same cost may consume a substantial fraction of the
per-frame time budget, reducing achievable throughput or violating
latency deadlines. The developer of an embedded computer-vision
application therefore faces a genuine dilemma: adopt the modular,
well-supported, ecosystem-integrated architecture that ROS~2 promotes, or
retain a monolithic single-process design that minimizes internal
communication at the price of maintainability and tooling. Resolving that
dilemma with evidence, rather than intuition, requires controlled
measurement on the class of hardware where the tension is sharpest.

\section{Problem Statement}
\label{sec:problem-statement}

A computer-vision application can be implemented as a single process in
which image acquisition, preprocessing, inference, and postprocessing
execute sequentially within one program. In this \emph{monolithic}
approach, data passes between processing stages as ordinary in-memory
function calls, with no serialization and no inter-process transport.
Alternatively, the same application can be decomposed into ROS~2 nodes
that communicate by exchanging messages over topics. This
\emph{node-based} architecture provides greater modularity, but it
introduces communication boundaries between processing stages, and each
boundary may add overhead in time and computational resources
\cite{maruyama2016, kronauer2021, kamiya2025, luo2025}. On an embedded
platform with limited CPU, memory, and energy, these additional costs may
affect the achievable throughput and latency of the application
\cite{kronauer2021, macenski2023}. The magnitude of the effect depends on
the workload, the degree of partitioning, and the hardware, and it
therefore requires experimental characterization rather than assumption.

This tension is not specific to robotics. In general software
engineering, the trade-off between monolithic and distributed
(microservice-style) architectures has been extensively studied, and the
evidence cautions against assuming that distribution is free: on a single
machine, a monolith has been shown to outperform its microservice-based
counterpart \cite{blinowski2022}, performance differences between the two
styles are measurable and workload-dependent \cite{tapia2020}, and a
recent critical review even identifies a trend of migrating back from
distributed to monolithic architectures, motivated in part by cost,
complexity, and performance \cite{felisberto2024}. A ROS~2 multi-node
system is, in this sense, a distributed architecture executed on a single
host, and ROS~2's own maintainers acknowledge the stakes: the composition
mechanism was created specifically to combat the performance issues of
the one-process-per-node model, and at least one commercial robot with
fewer than 60~MB of available RAM could not have shipped its ROS~2
application as separate processes at all \cite{macenski2023}. Studies of
full autonomous-driving stacks decompose end-to-end latency into
communication, computation, and idle components and identify latencies as
critical issues in production-scale ROS~2 systems \cite{betz2023a}, and
architectural variants such as containerized microservice deployments have
been evaluated precisely because these communication costs matter
\cite{betz2024}.

What the reviewed literature does not yet provide, however, is a
controlled answer at the scale of an embedded computer-vision pipeline.
Existing measurement studies characterize the middleware in isolation,
using synthetic payloads and publisher--subscriber chains
\cite{maruyama2016, kronauer2021, kamiya2025, luo2025, zhang2024}; studies
of mitigations such as node composition evaluate alternatives to the
multi-process baseline rather than the baseline decision itself
\cite{macenski2023}; and full-stack evaluations such as those of Autoware
conflate countless architectural and algorithmic factors
\cite{betz2023a, betz2024}. To the best of the author's knowledge, no
study in the reviewed literature holds the computer-vision workload,
model, input data, and hardware constant while varying only the
application architecture---monolithic process versus ROS~2 node
pipeline---and the degree of pipeline partitioning, on a Jetson-class
embedded platform and with the rclpy client library that many robotics
developers actually use.
% TODO: support the claim that Python/rclpy is widely used among ROS 2
% developers with adoption data (e.g., from the Portugal et al. survey's
% full text or the ROS 2 documentation); no source in the library states
% this explicitly.
The problem addressed by this thesis is precisely that gap: to quantify
the performance trade-off that accompanies the adoption of a node-based
architecture for a representative computer-vision workload on an embedded
platform, under controlled and reproducible conditions.

The relevance of this problem is threefold. First, it is timely: with
ROS~1 reaching end of life, teams are migrating perception systems to
ROS~2 now and must make architecture decisions with lasting consequences
\cite{portugal2025}. Second, it is consequential: the target hardware
class---exemplified by the Jetson Orin Nano---is a mainstream choice for
edge computer vision \cite{scalcon2024, alqahtani2026}, and the payloads
involved, image streams, are among the most demanding ones a ROS~2
pipeline routinely carries \cite{kamiya2025, macenski2023}. Third, it is
methodologically needed: the robotics community itself calls for
standardized, reproducible benchmarking of robotics computing
\cite{mayoral-vilches2024}, and workload-controlled comparisons of
architectural alternatives are a necessary ingredient of that effort. The
contribution of this thesis is an experimental characterization of the
performance trade-offs associated with implementing a computer-vision
pipeline as a set of ROS~2 nodes rather than as a monolithic process on an
embedded NVIDIA platform, providing practical guidance for developers who
must decide whether the modularity of ROS~2 carries a significant
performance penalty for this class of workload. The thesis does not
dispute the modularity benefits of ROS~2, nor does it attempt to establish
a universal performance ranking of robotics middleware.

\section{Research Questions}
\label{sec:research-questions}

The primary research question of this thesis is:

\begin{quote}
\emph{What is the performance impact of partitioning a computer-vision
pipeline into ROS~2 rclpy nodes compared with an equivalent single-process
Python implementation when executed on an NVIDIA Jetson Orin Nano?}
\end{quote}

In addition, a secondary research question is investigated:

\begin{quote}
\emph{How does increasing the degree of partitioning of the
computer-vision pipeline into ROS~2 nodes affect latency, throughput, and
resource utilization?}
\end{quote}

The study is specifically concerned with the performance trade-off
associated with introducing ROS~2 communication boundaries into an
otherwise equivalent computer-vision application. It is not concerned with
identifying the fastest possible way to run computer vision with ROS~2:
that question involves middleware implementations, QoS configurations,
deployment environments, and programming languages that are deliberately
held constant here, as described below.

\section{Objectives}
\label{sec:objectives}

\subsection{General Objective}

To experimentally characterize the performance difference between a
monolithic Python computer-vision pipeline and equivalent ROS~2 rclpy
implementations on an NVIDIA Jetson Orin Nano.

\subsection{Specific Objectives}

\begin{itemize}
  \item Develop a representative computer-vision pipeline capable of
        running on the Jetson Orin Nano.
  \item Implement the pipeline as a monolithic single-process Python
        application.
  \item Implement equivalent processing stages as ROS~2 rclpy nodes.
  \item Develop multiple ROS~2 configurations with different degrees of
        pipeline partitioning.
  \item Establish a reproducible benchmarking methodology.
  \item Measure end-to-end latency and processing throughput.
  \item Measure CPU, RAM, and GPU utilization.
  \item Optionally measure power consumption if reliable measurement is
        available.
  \item Compare the results statistically across configurations.
  \item Determine how the performance impact changes as the ROS~2
        pipeline becomes more highly partitioned.
\end{itemize}

\section{Scope and Limitations}
\label{sec:scope}

This thesis compares two implementations of the same computer-vision
application: a monolithic single-process Python implementation and a
multi-process ROS~2 implementation using rclpy. Both implementations share
the same programming language and runtime so that the programming language
is a controlled variable rather than a factor under investigation. The
underlying computer-vision algorithm, model, input data, hardware, and
relevant processing parameters remain equivalent between implementations.
The primary independent variable is the application architecture; a
secondary independent variable is the degree of ROS~2 partitioning. The
dependent variables are end-to-end latency, throughput, CPU, RAM, and GPU
utilization, and, optionally, power consumption.

A single ROS~2 distribution, a single DDS implementation, and a single
fixed QoS configuration are used, with a non-containerized deployment.
These choices are treated as controlled experimental conditions rather
than variables under investigation.

The following topics are explicitly outside the scope of this thesis:

\begin{itemize}
  \item Comparison of different programming languages
  \item C++ / rclcpp benchmarking
  \item ROS~2 composable nodes
  \item Comparison of different DDS implementations
  \item Comparison of different ROS~2 middleware configurations
  \item Systematic comparison of QoS policies
  \item Zenoh benchmarking
  \item Docker/containerization benchmarking
  \item Comparison of different ROS~2 distributions
  \item Comparison of different Jetson hardware platforms
  \item Comparison with other robotics middleware
  \item Comprehensive benchmarking of ROS~2 as a whole
  \item Evaluation of multiple unrelated computer-vision applications
\end{itemize}

These exclusions are deliberate. The thesis is not intended to answer what
the fastest possible way to run computer vision with ROS~2 is, but rather
what performance cost is associated with partitioning an equivalent
computer-vision application into ROS~2 nodes on an NVIDIA Jetson Orin
Nano. The excluded topics are discussed as potential future work in
Chapter~6. Additional limitations of the study, including the use of a
single hardware platform, a single computer-vision workload, a fixed DDS
and QoS configuration, and Python implementations, are scope limitations
of a controlled study rather than deficiencies, and are discussed together
with the conclusions.

\section{Thesis Organization}
\label{sec:organization}

The remainder of this thesis is organized as follows.

\textbf{Chapter~2 --- Background and Related Work} introduces computer
vision on embedded platforms, the architecture of the NVIDIA Jetson
platform, the ROS~2 architecture and its communication mechanisms, the
selected computer-vision workload, and the performance metrics used
throughout the thesis. It then reviews related work on ROS~2 benchmarks,
robotics middleware studies, embedded computer-vision benchmarks, and
studies of ROS~2 communication overhead, and identifies the research gap
addressed by this thesis.

\textbf{Chapter~3 --- Methodology} describes the experimental platform and
software environment, the computer-vision pipeline and its monolithic and
ROS~2 implementations, the experimental configurations, the benchmarking
procedure, the definition of each performance metric, the data-analysis
methods, and the measures taken to ensure reproducibility.

\textbf{Chapter~4 --- Results} presents the experimental results, including
throughput, end-to-end latency, CPU, memory, and GPU utilization, and, if
measured, power consumption, together with the comparison of the ROS~2
configurations with different degrees of partitioning.

\textbf{Chapter~5 --- Discussion} interprets the results, evaluates the
hypotheses, compares the findings with previous research, and discusses
unexpected observations and their possible explanations.

\textbf{Chapter~6 --- Conclusions and Future Work} answers the research
questions, summarizes the main findings, discusses the limitations of the
study, and outlines directions for future research.
